\begin{env}[5.10.11]
\label{IV.2.s5.10.11-fr}
Sean $X$ un preesquema localmente noetheriano, $Z$ una parte de $X$ estable
por especialización; se ha visto que
$\mathcal A=\mathcal H^0_{X/Z}(\mathcal O_X)$ es una
$\mathcal O_X$-álgebra cuasicoherente (5.9.3); el $X$-esquema
$X'=\operatorname{Spec}(\mathcal H^0_{X/Z}(\mathcal O_X))$ (\textbf{II}, 1.3.1) se
llama la $Z$-\emph{clausura} de $X$. Además, para todo $\mathcal O_X$-módulo
$\mathcal F$, $\mathcal H^0_{X/Z}(\mathcal F)$ es un $\mathcal A$-módulo,
que es cuasicoherente si $\mathcal F$ es un $\mathcal O_X$-módulo
cuasicoherente; en este último caso, existe por tanto un único
$\mathcal O_{X'}$-módulo $\mathcal F'$ tal que, designando por
$g:X'\to X$ el morfismo estructural (\textbf{II}, 1.4.3), se tenga
\begin{equation}
\mathcal H^0_{X/Z}(\mathcal F)=g_*(\mathcal F').
\tag{5.10.11.1}\label{IV.2.s5.10.11.1-fr}
\end{equation}
\end{env}

\begin{proposition}[5.10.12]
\label{IV.2.s5.10.12-fr}
Con las notaciones de \textup{(5.10.11)}:
\begin{enumerate}
\item[\textup{(i)}] Sea $x$ un punto de $X$. Para que el morfismo
\[
X'\times_X X_x\longrightarrow X_x(=\operatorname{Spec}(\mathcal O_{X,x}))
\]
obtenido de $g$ por localización sea un isomorfismo, es necesario y suficiente
que $\mathcal O_X$ sea $Z$-cerrado en el punto $x$ (lo que ocurre para todo
$x\in X\setmin Z$).
\item[\textup{(ii)}] Pongamos $Z'=g^{-1}(Z)$, y supongamos que $X'$ es localmente
noetheriano. Entonces $\mathcal F'$ es $Z'$-cerrado; si además $\mathcal F'$ es
un $\mathcal O_{X'}$-módulo coherente, se tiene
$\operatorname{prof}_{Z'}(\mathcal F')\geq2$.
\item[\textup{(iii)}] Supongamos que $\mathcal H^0_{X/Z}(\mathcal O_X)$ y
$\mathcal H^0_{X/Z}(\mathcal F)$ son coherentes. Entonces el morfismo
$g:X'\to X$ es finito; el conjunto $U$ de los $x\in X$ tales que se tenga
$\operatorname{prof}_{Z_x}(\widetilde{\mathcal O}_x)\geq2$ y
$\operatorname{prof}_{Z_x}(\widetilde{\mathcal F}_x)\geq2$ es abierto en
$X$ y se tiene $X\setmin U\subset Z$; además, $U$ es el mayor conjunto abierto de $X$
tal que la restricción $g^{-1}(U)\to U$ de $g$ sea un isomorfismo y que la
restricción $\mathcal F|U\to\mathcal F'|g^{-1}(U)$ del $g$-morfismo canónico
$\mathcal F\to\mathcal F'$ sea un isomorfismo.
\end{enumerate}
\end{proposition}

La afirmación (i) resulta de las definiciones, y (iii) es una consecuencia
inmediata de (5.10.10, (iii)). Para demostrar (ii), consideremos un abierto
$V$ de $X$ que contiene $Z\setmin X$, y su imagen inversa $V'=g^{-1}(V)$; si
$i:V\to X$ e $i':V'\to X'$ son las inyecciones

\oldpage[IV-2]{119}

canónicas, el homomorfismo canónico
$\rho_{X'/Z'}:\mathcal F'\to i'_*(\mathcal F'|V')$ es tal que
$g_*(\rho_{X'/Z'})$ sea el homomorfismo canónico
$\rho_{X/Z}:\mathcal H^0_{X/Z}(\mathcal F)
\to i_*(\mathcal H^0_{X/Z}(\mathcal F)|V)$
(\textbf{II}, 1.4.2), teniendo en cuenta (5.10.11.1). Como
$\mathcal H^0_{X/Z}(\mathcal F)$ es $Z$-cerrado (5.9.11), $\rho_{X/Z}$ es un
isomorfismo, por lo que ocurre lo mismo con $\rho_{X'/Z'}$. Como $X'\setmin Z'$ es
intersección de la familia filtrante de los $V'_\alpha=g^{-1}(V_\alpha)$,
donde $V_\alpha$ recorre la familia filtrante de los abiertos que contienen
$X\setmin Z$, se deduce que $\mathcal F'$ es $Z'$-cerrado cuando $X'$ es localmente
noetheriano, en virtud de (5.9.1).

\begin{env}[5.10.13]
\label{IV.2.s5.10.13-fr}
Vamos a aplicar ahora los resultados precedentes al caso en que $Z$ es uno de
los conjuntos $Z^{(n)}(X)$ (o simplemente $Z^{(n)}$), definido como el conjunto
de los $x\in X$ tales que $\dim(\mathcal O_x)\geq n$; está claro que
$Z^{(n)}$ es estable por especialización; para que una parte cerrada $T$ de $X$
esté contenida en $Z^{(n)}$, es necesario y suficiente que
$\operatorname{codim}(T,X)\geq n$. Nos interesaremos aquí por el caso $n=2$.
\end{env}
